\begin{abstract}
Can more computation compensate for an unavailable denoising scale?
We study this question for a fixed RED gradient update initialized
by backprojection. On a positive-width family of three-template
Gaussian mixtures, every controller confined to two variance ranges
has positive limiting recovery risk, regardless of its history
dependence or finite budget. An intermediate band instead permits
entry into the correct terminal region in at most 25 denoiser calls,
followed by logarithmic refinement to vanishing risk.
The mechanism is competition between templates: low and high scales
preserve a common wrong region, while intermediate scales suppress
a competing template enough to open a recovery channel.
For two templates aligned with a sensing eigenvector, endpoint
scales already match any strict crossing up to the midpoint.
Geometric tests connect local drift, full trajectories, and recovery
risk, and select bands and entrance budgets before trajectory
evaluation. Frozen experiments test these predictions and compare
independently tuned shared and observation-gated platforms.
Photo-patch experiments delineate a complementary regime where
terminal refinement erases schedule differences.
\end{abstract}
